\begin{table}[htbp]
  \centering
  \caption{Testing whether wage gaps relative to 11--19 workers changed over time}
  \label{tab:firm-size-premium-trend-test-11-19-reference}
  \small
  \begin{tabular}{lcccc}
    \toprule
    Firm size & Annual trend & S.E. & $p$-value & Implied 2008--2025 change \\
    \midrule
    Solo & -0.0027 & (0.0039) & 0.491 & -4.5\% \\
    2-3 & 0.0005 & (0.0016) & 0.746 & 0.9\% \\
    4-5 & 0.0005 & (0.0014) & 0.736 & 0.8\% \\
    6-10 & -0.0006 & (0.0011) & 0.611 & -1.0\% \\
    20-30 & -0.0009 & (0.0010) & 0.392 & -1.5\% \\
    31-50 & 0.0000 & (0.0012) & 0.988 & 0.0\% \\
    51-100 & -0.0001 & (0.0014) & 0.940 & -0.2\% \\
    101+ & 0.0017 & (0.0012) & 0.175 & 3.0\% \\
    \bottomrule
  \end{tabular}
  \vspace{0.3em}
  \begin{minipage}{0.95\textwidth}
  \footnotesize
  Notes: The annual trend is the coefficient on the interaction between firm-size category and a linear year trend, with workers in firms with 11--19 workers as the omitted category. The specification controls for gender, age, age squared, education, formality, and sector-year fixed effects, and uses GEIH expansion weights. Standard errors are clustered by sector. The $p$-value corresponds to the two-sided test that the trend differs from zero. A positive trend means that the listed category gained ground relative to 11--19 workers; a negative trend means that it lost ground. The final column reports $100\times[\exp(17\hat{\delta})-1]$, the implied change in the wage ratio between the listed category and 11--19 workers between 2008 and 2025.
  \end{minipage}
\end{table}
